% ============================================================
% Section 01: 固体物理简答题汇总（一）——晶体结构、结合与缺陷
% ============================================================
\section{固体理论：简答题汇总（一）——晶体结构、结合与缺陷}
\label{sec:solid-state-short-1}

\begin{tipbox}{关键词标签}
晶体结构 · 布拉维格子 · 密勒指数 · 倒格子 · 晶体结合 · 缺陷 · 位错 · 致密度 · 配位数
\end{tipbox}

\begin{insightbox}{本章定位}
本章汇总固体物理中关于晶体结构、晶体结合、晶格缺陷等基础概念的简答题。这些问题是固体物理考试的常考点，注重概念理解的准确性与表述的完整性。
\end{insightbox}

\vspace{0.5em}

\begin{enumerate}[label=\textbf{\arabic*.}, leftmargin=2em, itemsep=0.8em]

% --------------------------------------------------------
\item \textbf{晶体和非晶体的区别是什么？}

晶体的原子排列具有严格的周期性，即长程有序（long-range order），并表现出以下特征：
\begin{itemize}[nosep]
    \item 有固定的熔点；
    \item 物理性质各向异性（anisotropy）；
    \item 具有自限性，能自发形成规则几何外形（解理性）。
\end{itemize}
非晶体则仅具有短程有序（short-range order），没有长程的平移对称性，没有固定熔点（存在玻璃转变温度），物理性质各向同性。

% --------------------------------------------------------
\item \textbf{请说明晶体分类及其特点。}

晶体按结合力类型可分为五类：
\begin{enumerate}[nosep, label=(\arabic*)]
    \item \textbf{离子晶体}：正负离子间通过库仑相互作用结合，如 NaCl、CsCl；键强较高，熔点高、硬度大、常温下为绝缘体。
    \item \textbf{共价晶体}：相邻原子靠共用电子对（共价键）结合，如金刚石、Si、Ge；具有方向性和饱和性，结合力强，硬度高。
    \item \textbf{分子晶体}：原子或分子靠范德瓦尔斯力（瞬时偶极--诱导偶极相互作用）结晶，如固态 Ar、CO$_2$（干冰）；结合力弱，熔点低、质地软。
    \item \textbf{金属晶体}：结合作用来源于价电子的共有化（自由电子气与离子实之间的库仑作用），如 Cu、Al；无方向性，具有良好的导电性、导热性和延展性。
    \item \textbf{氢键晶体}：氢原子同时与其他两个电负性较强的原子形成氢键结合，如冰、HF；键强介于共价键与范德瓦尔斯力之间。
\end{enumerate}

% --------------------------------------------------------
\item \textbf{在晶体衍射中，为什么不能用可见光？}

晶体中原子间距的量级为 \AA\ ($10^{-10}\,$m)，而可见光波长约为 $400$--$700\,$nm，即 $4000$--$7000$\AA。根据布拉格衍射条件
\[
    2d\sin\theta = n\lambda,
\]
要产生可观测的衍射，波长 $\lambda$ 必须与晶面间距 $d$ 相当（$\lambda\sim d$）。可见光波长远大于晶面间距，因此无法满足布拉格条件，不能用于晶体衍射。晶体衍射需要使用波长与原子间距相当的 \textbf{X 射线}（$\lambda\sim 1$\AA）、\textbf{电子束}（德布罗意波长可调）或 \textbf{中子束}（对轻元素和磁性敏感）。

% --------------------------------------------------------
\item \textbf{晶体的结合能、内能、原子间的相互作用势能有何区别？}

\begin{itemize}[nosep]
    \item \textbf{结合能（内聚能，Cohesive Energy）}：与分离成各个孤立原子相比，各原子聚合形成晶体后系统能量的\textbf{下降值}，即把晶体全部分离成孤立原子所需要的能量。
    \item \textbf{内能（Internal Energy）}：晶体内部一切微观粒子的一切运动形式所具有的能量总和，包括电子动能、晶格振动能、零点能等。
    \item \textbf{原子间的相互作用势能}：两个原子之间由于相互作用而产生的势能，通常可写成吸引势与排斥势之和，如Lennard-Jones势
    \[
        \varphi(r) = 4\varepsilon\Bigl[\Bigl(\frac{\sigma}{r}\Bigr)^{12} - \Bigl(\frac{\sigma}{r}\Bigr)^{6}\Bigr].
    \]
\end{itemize}
三者的关系：晶体总内能 $=$ 所有原子对相互作用势能之和 $+$ 各粒子动能 $+$ 其他贡献；结合能 $=$ 孤立原子总能量 $-$ 晶体内能（取绝对值）。

% --------------------------------------------------------
\item \textbf{晶体具有哪些共性？}

晶体具有以下共性：
\begin{enumerate}[nosep, label=(\arabic*)]
    \item \textbf{长程有序}：原子排列具有三维周期性；
    \item \textbf{自限性}：能自发形成规则几何外形；
    \item \textbf{各向异性}：物理性质（如热膨胀系数、电导率、折射率等）随方向不同而异；
    \item \textbf{对称性}：在旋转、反演、镜面反映等对称操作下保持不变；
    \item \textbf{固定的熔点}：晶体在确定温度下发生固--液相变；
    \item \textbf{解理性}：晶体常沿某些特定晶面（解理面）容易劈裂。
\end{enumerate}

% --------------------------------------------------------
\item \textbf{什么是晶格，什么是格点，什么是点阵？}

\begin{itemize}[nosep]
    \item \textbf{格点（Lattice Point）}：将基元抽象为几何点，代表原子平衡位置的几何点；
    \item \textbf{点阵（Lattice）}：格点的集合，是一种数学抽象，反映晶体粒子排列的周期性；
    \item \textbf{晶格（Crystal Lattice）}：连接格点后抽象出的几何网格。
\end{itemize}
三者关系为：格点 $\to$ 点阵（格点的集合）$\to$ 晶格（点阵的几何表示），而实际晶体结构 $=$ 晶格点阵 $+$ 基元（basis）。

% --------------------------------------------------------
\item \textbf{请说明晶体的周期性和对称性特征。}

\textbf{周期性（平移对称性）}：晶体中原子或离子在空间按一定规律作周期性的重复排列，可用平移矢量
\[
    \mathbf{R}_n = n_1\mathbf{a}_1 + n_2\mathbf{a}_2 + n_3\mathbf{a}_3
\]
来描述，其中 $\mathbf{a}_1, \mathbf{a}_2, \mathbf{a}_3$ 为基矢，$n_1, n_2, n_3$ 为任意整数。

\textbf{对称性}：晶体在旋转、反演、镜面反映、平移等对称操作下保持不变的性质。晶体的宏观对称性受平移对称性限制，只可能有 $1, 2, 3, 4, 6$ 次旋转轴（不存在5次或7次以上旋转轴）。

% --------------------------------------------------------
\item \textbf{什么是密勒指数？}

以结晶学原胞基矢 $\mathbf{a}, \mathbf{b}, \mathbf{c}$ 为坐标轴，选任一结点为原点，求出晶面族中离原点最近的第一个晶面在坐标轴上的截距（分别记为 $a/h, b/k, c/l$），取倒数并化为互质整数 $(h\,k\,l)$，即为该晶面的密勒指数（Miller Indices）。

\textbf{特点}：密勒指数与晶面法线方向密切相关；晶面间距 $d_{hkl}$ 可通过倒格子基矢计算。负截距用上划线表示，如 $(\bar{1}11)$。

% --------------------------------------------------------
\item \textbf{请说明离子性、共价性、金属性和范德瓦尔斯性结合力的特点。}

\begin{itemize}[nosep]
    \item \textbf{离子性结合}：由正负离子间的库仑吸引力产生，方向性不显著，结合力较强，晶体熔点高、硬度大、常温绝缘（高温离子导电）。
    \item \textbf{共价性结合}：相邻原子通过共用电子对形成共价键，具有\textbf{方向性}和\textbf{饱和性}，结合力很强，晶体硬度高、熔点高。
    \item \textbf{金属性结合}：来源于价电子的共有化，即自由电子气与离子实之间的库仑相互作用，无方向性和饱和性，具有良好的导电性、导热性和延展性。
    \item \textbf{范德瓦尔斯性结合}：由分子间的瞬时偶极--诱导偶极相互作用产生，结合力很弱，晶体熔点低、质地软。
\end{itemize}

% --------------------------------------------------------
\item \textbf{什么是倒格子？}

设一晶格的基矢为 $\mathbf{a}_1, \mathbf{a}_2, \mathbf{a}_3$，若另一格子的基矢 $\mathbf{b}_1, \mathbf{b}_2, \mathbf{b}_3$ 满足
\[
    \mathbf{a}_i \cdot \mathbf{b}_j = 2\pi\delta_{ij},
\]
则以 $\mathbf{b}_1, \mathbf{b}_2, \mathbf{b}_3$ 为基矢的格子称为原晶格的\textbf{倒格子}（Reciprocal Lattice）。倒格子所在的空间称为倒空间，是晶体衍射理论（X射线衍射、电子衍射）和能带理论的基础。

倒格子基矢的显式表达式为
\[
    \mathbf{b}_1 = \frac{2\pi(\mathbf{a}_2\times\mathbf{a}_3)}{\mathbf{a}_1\cdot(\mathbf{a}_2\times\mathbf{a}_3)},\quad
    \mathbf{b}_2 = \frac{2\pi(\mathbf{a}_3\times\mathbf{a}_1)}{\mathbf{a}_1\cdot(\mathbf{a}_2\times\mathbf{a}_3)},\quad
    \mathbf{b}_3 = \frac{2\pi(\mathbf{a}_1\times\mathbf{a}_2)}{\mathbf{a}_1\cdot(\mathbf{a}_2\times\mathbf{a}_3)}.
\]

% --------------------------------------------------------
\item \textbf{什么是致密度？}

致密度（Packing Fraction 或 Atomic Packing Factor, APF）是指晶胞内原子所占体积与晶胞总体积之比，用于衡量原子堆积的紧密程度。

\textbf{常见结构的致密度}：
\begin{center}
\renewcommand{\arraystretch}{1.4}
\begin{tabular}{|l|c|c|c|}
\hline
\textbf{结构} & \textbf{配位数} & \textbf{晶胞原子数} & \textbf{致密度} \\
\hline
简单立方 (sc) & 6 & 1 & $\pi/6 \approx 0.52$ \\
体心立方 (bcc) & 8 & 2 & $\pi\sqrt{3}/8 \approx 0.68$ \\
面心立方 (fcc) & 12 & 4 & $\pi\sqrt{2}/6 \approx 0.74$ \\
六方密堆积 (hcp) & 12 & 6 & $\pi\sqrt{2}/6 \approx 0.74$ \\
金刚石结构 & 4 & 8 & $\pi\sqrt{3}/16 \approx 0.34$ \\
\hline
\end{tabular}
\end{center}

% --------------------------------------------------------
\item \textbf{请说明原子负电性与晶体结构的关系。}

原子的负电性（电负性，Electronegativity）表示原子吸引电子的能力。鲍林（Pauling）将电负性定义为
\[
    \chi = 0.18\,(E_{\text{ion}} + E_{\text{aff}}),
\]
其中 $E_{\text{ion}}$ 为电离能，$E_{\text{aff}}$ 为电子亲和能。

\textbf{与晶体结构的关系}：
\begin{itemize}[nosep]
    \item 两种原子的负电性差越大，越趋向于形成\textbf{离子晶体}（如 NaCl，$\Delta\chi \approx 2.23$）；
    \item 负电性差较小则趋向于形成\textbf{共价晶体}（如 GaAs，$\Delta\chi \approx 0.37$）；
    \item 同种原子或电负性相近的金属原子形成\textbf{金属晶体}。
\end{itemize}

% --------------------------------------------------------
\item \textbf{请说明原胞和晶胞的特点和区别。}

\begin{itemize}[nosep]
    \item \textbf{原胞（Primitive Cell，固体物理学原胞）}：只反映晶格周期性的\textbf{最小}重复单元，结点只在顶点上，内部和面上不含其他结点。原胞是体积最小的周期性结构单元，每个原胞只含一个格点。
    \item \textbf{晶胞（Conventional Cell / 结晶学原胞 / 单胞）}：同时计及晶格的周期性和\textbf{对称性}的重复单元，结点既可在顶角上，也可在体心或面心处。晶胞通常比原胞大，但能更好地体现晶体的宏观对称性。
\end{itemize}

\textbf{举例}：体心立方（bcc）的晶胞是立方体（含2个原子），而原胞是菱方体（含1个原子）；面心立方（fcc）的晶胞含4个原子，原胞含1个原子。

% --------------------------------------------------------
\item \textbf{什么是配位数？}

配位数（Coordination Number）是指晶体中每个原子周围\textbf{最近邻}原子的数目。

\textbf{常见结构的配位数}：
\begin{itemize}[nosep]
    \item 简单立方（sc）：6
    \item 体心立方（bcc）：8
    \item 面心立方（fcc）：12
    \item 六方密堆积（hcp）：12
    \item 金刚石结构：4
    \item NaCl 结构：6
    \item CsCl 结构：8
\end{itemize}

% --------------------------------------------------------
\item \textbf{请说明布拉维格子和复式格子的特点和区别。}

\begin{itemize}[nosep]
    \item \textbf{布拉维格子（Bravais Lattice，简单格子）}：由矢量 $\mathbf{R}_n = n_1\mathbf{a}_1 + n_2\mathbf{a}_2 + n_3\mathbf{a}_3$ 全部端点的集合构成，每个基元中只有一个原子（或一个等价点），所有格点在位置上\textbf{完全等价}。三维空间中共有14种布拉维格子。
    \item \textbf{复式格子（Complex Lattice）}：每个基元中包含一个以上的原子或离子，可以看作若干相同的布拉维格子相互套构而成。
\end{itemize}

\textbf{举例}：体心立方结构的碱金属和面心立方结构的 Au、Ag、Cu 是简单格子；NaCl、CsCl、金刚石、ZnS 是复式格子。

% --------------------------------------------------------
\item \textbf{为什么许多金属为密积结构？}

金属键不具有方向性和饱和性，金属原子可看作刚性球。为使系统能量最低（结合能最大），原子趋向于尽可能紧密的排列方式。面心立方（fcc）和六方密堆积（hcp）是两种最密的堆积方式（致密度均为 $0.74$，配位数为12），因此许多金属采取这两种密积结构。

\end{enumerate}

% ============================================================
% Section 02: 固体物理简答题汇总（二）——晶格振动、热学性质与能带电子
% ============================================================
\section{固体理论：简答题汇总（二）——晶格振动、热学性质与能带电子}
\label{sec:solid-state-short-2}

\begin{tipbox}{关键词标签}
德拜模型 · 声子 · 电子热容 · 能带理论 · 布洛赫定理 · 紧束缚近似 · 有效质量 · 自由电子模型
\end{tipbox}

\begin{insightbox}{本章定位}
本章汇总固体物理中关于晶格振动、固体热容、金属电子论及能带理论的简答题。这些问题涉及从经典到量子理论的过渡，是理解固体物理核心框架的关键。
\end{insightbox}

\vspace{0.5em}

\begin{enumerate}[label=\textbf{\arabic*.}, leftmargin=2em, itemsep=0.8em, resume]

% --------------------------------------------------------
\item \textbf{在低温下，德拜模型为什么与实验相符？}

低温时，对晶体热容的贡献主要来自\textbf{长波声学模}（long-wavelength acoustic modes）。德拜模型假设以连续介质的弹性波代表格波，恰好在长波极限（$q\to 0$）下与真实的声学支色散关系一致。低温下只有能量 $\hbar\omega < k_{\text{B}}T$ 的声子模式被热激发，而这些模式正好处于长波区域，德拜模型的线性色散 $\omega = v_s q$ 很好地近似了实际情况。因此德拜模型在低温下与实验吻合良好，热容按 $T^3$ 律变化：
\[
    C_V = \frac{12\pi^4}{5} N k_{\text{B}} \Bigl(\frac{T}{\Theta_D}\Bigr)^3.
\]

% --------------------------------------------------------
\item \textbf{请解释金属、半导体和绝缘体的导电性原理。}

\begin{itemize}[nosep]
    \item \textbf{导体（金属）}：除去完全充满的一系列能带（满带）外，还有只是\textbf{部分被电子填充}的能带（导带）。导带中的电子在外电场作用下可获得净动量，从而导电。
    \item \textbf{绝缘体}：电子恰好填满最低的一系列能带（满带），再高的能带全空。满带中电子的 $k$ 态分布对称，不产生净电流，因此不导电。禁带宽度 $E_g \gtrsim 3\,$eV。
    \item \textbf{半导体}：能带填充情况介于导体和绝缘体之间。禁带宽度较小（$E_g \sim 0.1$--$2\,$eV），通过热激发可使导带中有少数电子，或满带中缺少少数电子（空穴），从而具备一定的导电性。本征半导体的载流子浓度随温度指数增长：$n_i \propto \exp(-E_g/2k_{\text{B}}T)$。
\end{itemize}

% --------------------------------------------------------
\item \textbf{金属淬火后为什么变硬？}

淬火（quenching）是将金属加热到高温后快速冷却的热处理工艺。金属淬火后变硬的主要原因包括：
\begin{enumerate}[nosep, label=(\arabic*)]
    \item \textbf{马氏体相变}（对钢而言）：快速冷却抑制了扩散型相变，奥氏体（fcc）通过无扩散切变转变为马氏体（体心四方，bct）。马氏体内部存在大量晶格畸变、高密度位错和显微孪晶，严重阻碍位错运动，因而硬度大幅提高。
    \item \textbf{过饱和固溶体}（对铝合金等）：快速冷却使高温下的固溶原子来不及析出，形成过饱和固溶体。后续时效处理可析出细小弥散的第二相粒子，通过\textbf{析出强化}（precipitation hardening）机制阻碍位错运动。
    \item \textbf{晶粒细化}：快速冷却可细化晶粒，晶界增多，通过\textbf{细晶强化}（Hall-Petch效应）提高硬度。
\end{enumerate}

% --------------------------------------------------------
\item \textbf{在位错滑移时，刃位错上原子受力和螺位错原子受力各有什么特点？}

\begin{itemize}[nosep]
    \item \textbf{刃位错（Edge Dislocation）}：存在多余的半原子面，位错线方向与滑移方向\textbf{垂直}。刃位错的滑移由\textbf{切应力}（shear stress）驱动，原子沿\textbf{垂直于位错线}的方向在滑移面内运动。刃位错只能在其唯一的滑移面内滑移，不能进行交滑移（cross-slip）。
    \item \textbf{螺位错（Screw Dislocation）}：原子排列呈螺旋状，位错线方向与滑移方向\textbf{平行}。螺位错的滑移同样由切应力驱动，但原子沿\textbf{平行于位错线}的方向运动。螺位错没有确定的滑移面，可在包含位错线的任意平面内滑移，因此可以发生交滑移。
\end{itemize}

% --------------------------------------------------------
\item \textbf{请说明金刚石的结构特点。}

金刚石结构属于\textbf{复式格子}，可看作由两个面心立方子晶格沿体对角线方向平移 $1/4$ 体对角线长度套构而成。

\textbf{结构参数}：
\begin{itemize}[nosep]
    \item 晶胞类型：面心立方（fcc）布拉维格子 + 双原子基元
    \item 每个晶胞含 8 个原子（8 个顶角 + 6 个面心 + 4 个内部）
    \item 每个原胞含 2 个原子
    \item 配位数：4（正四面体配位）
    \item 致密度：$\pi\sqrt{3}/16 \approx 0.34$
\end{itemize}

\textbf{键合方式}：$\text{sp}^3$ 杂化的共价键，键角 $109.5°$，结合力非常强，是自然界中已知最硬的物质之一。金刚石是绝缘体，禁带宽度约 $5.5\,$eV。

% --------------------------------------------------------
\item \textbf{量子理论中如何解释金属电子对固体热容的贡献？}

根据量子自由电子理论，金属中的电子服从\textbf{费米-狄拉克分布}。在温度 $T > 0$ 时，只有费米能级 $E_F$ 附近约 $k_{\text{B}}T$ 能量范围内的电子能够被热激发到更高能态（费米面附近的"热激发层"）。

这部分电子数约占电子总数的 $k_{\text{B}}T/E_F$，因此电子热容与温度 $T$ 成正比：
\[
    C_e = \gamma T, \qquad \gamma = \frac{\pi^2}{3} k_{\text{B}}^2 g(E_F),
\]
其中 $g(E_F)$ 为费米能级处的态密度。该结果成功解释了经典德鲁德模型无法说明的实验事实：常温下电子热容远小于晶格热容。

% --------------------------------------------------------
\item \textbf{为什么常温下电子对热容的贡献很小？}

常温下，$k_{\text{B}}T \approx 0.025\,$eV，而金属的费米能 $E_F$ 约为 $1$--$10\,$eV（如铜的 $E_F \approx 7\,$eV）。因此
\[
    \frac{k_{\text{B}}T}{E_F} \sim 10^{-2}\text{--}10^{-3} \ll 1.
\]

由泡利不相容原理，费米球内部深处的电子态已被填满，无法被热激发；只有费米面附近极少量（约 $k_{\text{B}}T/E_F$ 比例）的电子能够参与热激发。因此常温下电子热容 $C_e \sim 10^{-2}Nk_{\text{B}}$，远小于晶格振动热容（杜隆-珀蒂值 $C_V = 3Nk_{\text{B}}$，德拜温度以上趋于此值），对总热容的贡献通常不到 $1\%$。

% --------------------------------------------------------
\item \textbf{什么是声子？}

声子（Phonon）是\textbf{格波的能量量子}。晶格振动的总能量可以量子化为
\[
    E = \sum_{q,s} \Bigl(n_{q,s} + \frac{1}{2}\Bigr)\hbar\omega_s(q),
\]
其中 $n_{q,s}$ 为波矢 $q$、支指标 $s$ 的声子数。声子能量为 $\hbar\omega$，准动量为 $\hbar q$。

\textbf{关键性质}：
\begin{itemize}[nosep]
    \item 声子不是真正的粒子，而是描述晶格振动量子化的\textbf{准粒子}；
    \item 声子是\textbf{玻色子}，遵从玻色-爱因斯坦统计；
    \item 声子数不守恒，可被激发和湮灭（与光子类似）；
    \item 声子之间的相互作用（通过非简谐项）是晶格热阻的来源。
\end{itemize}

% --------------------------------------------------------
\item \textbf{当电子的波矢落在布里渊区边界上时，为什么其有效质量与真实质量有显著差别？}

在布里渊区边界上，电子波满足\textbf{布拉格反射条件} $2k\cdot G = G^2$，电子波受到晶格周期势场的强烈散射，入射波与反射波相干叠加形成驻波。此时能带变得平坦，能量对波矢的二阶导数 $|d^2E/dk^2|$ 显著减小。

电子的有效质量定义为
\[
    m^* = \hbar^2 \Bigl(\frac{\partial^2 E}{\partial k^2}\Bigr)^{-1},
\]
因此在布里渊区边界附近 $m^*$ 可能与自由电子质量 $m_e$ 相差很大。在能带底 $m^* > 0$，在能带顶 $m^* < 0$（负有效质量描述的是空穴行为）。

% --------------------------------------------------------
\item \textbf{什么是非简谐效应？}

在晶格振动势能展开中，仅保留位移的二次项（\textbf{简谐近似}）可导出独立的声子图像，成功解释晶格比热容（爱因斯坦模型、德拜模型），但\textbf{无法解释热膨胀和热传导}等现象。

当考虑位移的三次及更高次项时，即为\textbf{非简谐效应}（anharmonic effect）。非简谐项导致：
\begin{itemize}[nosep]
    \item 声子之间存在相互作用（声子散射），是晶格热阻的来源；
    \item 热膨胀（格林爱森参数 $\gamma$ 描述）；
    \item 声子倒逆过程（Umklapp process，U过程），对热导率有重要影响。
\end{itemize}

% --------------------------------------------------------
\item \textbf{请说明研究金属热容的意义。}

研究金属热容具有以下重要意义：
\begin{enumerate}[nosep, label=(\arabic*)]
    \item \textbf{电子结构信息}：通过低温电子热容项 $\gamma T$ 的实验测定，可提取费米能级处的态密度 $g(E_F)$，进而了解金属的电子能带结构。
    \item \textbf{验证理论模型}：检验德拜模型、爱因斯坦模型及自由电子气理论的适用范围和精度。
    \item \textbf{获取特征参数}：测定德拜温度 $\Theta_D$、电子-声子耦合常数等重要物理量。
    \item \textbf{研究相变}：通过热容的异常（跃变或发散）探测超导转变、磁性相变、结构相变等。
    \item \textbf{超导电性}：BCS理论预言低温下超导态的电子热容呈指数激活行为 $C_{es} \propto \exp(-\Delta/k_{\text{B}}T)$，与正常态的线性行为截然不同。
\end{enumerate}

% --------------------------------------------------------
\item \textbf{请说明金属接触电势差如何形成。}

两种不同金属相互接触时，由于两者的\textbf{费米能级不同}（取决于各自的功函数 $\phi$），电子会从费米能级较高的金属（功函数较小）流向费米能级较低的金属（功函数较大），直到两者的费米能级在界面处拉平（热力学平衡条件）。

结果在接触面两侧形成\textbf{电势差}，即接触电势差（contact potential difference）：
\[
    V_A - V_B = \frac{\phi_B - \phi_A}{e}.
\]
该电势差阻止了电子的进一步净流动，使两金属内部费米能级相等。

% --------------------------------------------------------
\item \textbf{什么是极化子？}

在\textbf{离子晶体}中，运动的电子会使周围晶格发生极化，正负离子产生相对位移，形成局域的极化畸变区域。电子与其引起的极化场一起运动的准粒子称为\textbf{极化子}（Polaron）。

\textbf{分类}：
\begin{itemize}[nosep]
    \item \textbf{大极化子}（Large Polaron）：极化区域远大于晶格常数，可用连续介质近似描述；
    \item \textbf{小极化子}（Small Polaron）：极化区域局限于单个晶胞附近，需要离散晶格模型。
\end{itemize}

极化子的有效质量比\textbf{裸电子}（bare electron）大，因为电子运动时还要拖着周围的极化畸变一起运动。极化子概念对理解离子晶体和氧化物中的输运性质非常重要。

% --------------------------------------------------------
\item \textbf{什么是布洛赫波？}

布洛赫波（Bloch Wave）是\textbf{周期性势场}中单电子薛定谔方程的解，其形式为
\[
    \psi_{\mathbf{k}}(\mathbf{r}) = e^{i\mathbf{k}\cdot\mathbf{r}} u_{\mathbf{k}}(\mathbf{r}),
\]
其中 $u_{\mathbf{k}}(\mathbf{r})$ 具有晶格的周期性：$u_{\mathbf{k}}(\mathbf{r} + \mathbf{R}_n) = u_{\mathbf{k}}(\mathbf{r})$。

\textbf{物理图像}：布洛赫波是一个被周期函数 $u_{\mathbf{k}}(\mathbf{r})$ \textbf{调幅}的平面波。它描述的是在晶格周期势场中运动的电子态，既不像自由电子那样完全自由，也不像束缚电子那样完全局域。布洛赫波是能带理论的基础，保证了晶体中电子可以无散射地传播（在理想周期势中电阻为零）。

% --------------------------------------------------------
\item \textbf{请说明布洛赫定理和布洛赫函数。}

\textbf{布洛赫定理}：在周期性势场 $V(\mathbf{r} + \mathbf{R}_n) = V(\mathbf{r})$ 中运动的电子，其定态薛定谔方程的解可写为
\[
    \psi_{\mathbf{k}}(\mathbf{r}) = e^{i\mathbf{k}\cdot\mathbf{r}} u_{\mathbf{k}}(\mathbf{r})
\]
的形式，其中 $u_{\mathbf{k}}(\mathbf{r} + \mathbf{R}_n) = u_{\mathbf{k}}(\mathbf{r})$ 具有晶格的周期性。等价表述：波函数满足
\[
    \psi_{\mathbf{k}}(\mathbf{r} + \mathbf{R}_n) = e^{i\mathbf{k}\cdot\mathbf{R}_n} \psi_{\mathbf{k}}(\mathbf{r}).
\]

\textbf{布洛赫函数}：即形如 $\psi_{\mathbf{k}}(\mathbf{r}) = e^{i\mathbf{k}\cdot\mathbf{r}} u_{\mathbf{k}}(\mathbf{r})$ 的波函数。布洛赫定理是能带理论的数学基础。

% --------------------------------------------------------
\item \textbf{请说明紧束缚近似模型的思想和主要结论。}

\textbf{思想}：紧束缚近似（Tight-Binding Approximation, TBA）认为电子在原子附近主要受该原子的势场束缚，不同原子间的相互作用可看作微扰。将晶体中的电子波函数用孤立原子轨道的线性组合（LCAO，Linear Combination of Atomic Orbitals）来表示：
\[
    \psi_{\mathbf{k}}(\mathbf{r}) = \frac{1}{\sqrt{N}} \sum_{\mathbf{R}_n} e^{i\mathbf{k}\cdot\mathbf{R}_n} \phi(\mathbf{r} - \mathbf{R}_n).
\]

\textbf{主要结论}：
\begin{enumerate}[nosep, label=(\arabic*)]
    \item 电子的能带由孤立原子的离散能级\textbf{展宽}而成；
    \item 能带宽度取决于最近邻原子间的\textbf{交叠积分}（跃迁积分 $J$），原子间距越小，轨道交叠越大，能带越宽；
    \item 对于简单立方晶格 s 态电子，色散关系为
    \[
        E(\mathbf{k}) = E_s - J_0 - 2J_1(\cos k_x a + \cos k_y a + \cos k_z a);
    \]
    \item 紧束缚近似适用于\textbf{内层电子}和\textbf{窄能带}体系（如过渡金属的 d 电子）。
\end{enumerate}

% --------------------------------------------------------
\item \textbf{请说明陶瓷中晶界对材料性能的影响。}

陶瓷是多晶体，由大量取向不同的晶粒和它们之间的晶界（grain boundary）组成。晶界对陶瓷性能的影响包括：
\begin{enumerate}[nosep, label=(\arabic*)]
    \item \textbf{力学性能}：晶界对位错运动有阻碍作用（Hall-Petch强化），细化晶粒可提高强度和韧性；但高温下晶界滑移会导致蠕变。
    \item \textbf{电学性能}：晶界是杂质偏析和缺陷集中的区域，形成势垒或导电通道。在功能陶瓷中，晶界的电学特性往往是器件功能的基础：
    \begin{itemize}[nosep]
        \item PTC热敏电阻：晶界的肖特基势垒随温度变化导致电阻突变；
        \item 压敏电阻（ZnO）：晶界势垒控制非线性伏安特性。
    \end{itemize}
    \item \textbf{热学性能}：晶界会散射热声子，降低热导率。
    \item \textbf{扩散行为}：晶界扩散系数远高于晶内扩散，影响烧结和离子导电。
\end{enumerate}

% --------------------------------------------------------
\item \textbf{请说明经典的自由电子理论的要点及应用。}

\textbf{要点}：德鲁德（Drude, 1900）提出。将金属中的价电子视为\textbf{经典的自由粒子}（理想气体），在离子实的均匀正电荷背景中自由运动，服从\textbf{经典统计}（麦克斯韦-玻尔兹曼分布）。电子与离子实之间的碰撞导致电阻，碰撞后电子速度随机化。

\textbf{成功之处}：
\begin{itemize}[nosep]
    \item 解释欧姆定律 $j = \sigma E$；
    \item 解释电导与热导的关系——维德曼-弗兰兹定律（Wiedemann-Franz Law）：$\kappa/\sigma T = L$（洛伦兹数，实验值约 $2.44\times 10^{-8}\,\text{W}\cdot\Omega/\text{K}^2$）。
\end{itemize}

\textbf{失败之处}：无法解释电子比热容（理论值比实验值大两个数量级）、磁化率（泡利顺磁性）及霍耳系数的反常等。

% --------------------------------------------------------
\item \textbf{请说明自由电子近似的理论要点。}

\textbf{要点}：索末菲（Sommerfeld, 1928）在德鲁德模型基础上进行量子化改造：
\begin{enumerate}[nosep, label=(\arabic*)]
    \item 电子服从\textbf{费米-狄拉克统计}，满足\textbf{泡利不相容原理}；
    \item 电子处于周期性边界条件（Born-von Karman边界条件）下的三维空间中，视为自由粒子，波函数为平面波 $\psi_{\mathbf{k}}(\mathbf{r}) = \frac{1}{\sqrt{V}}e^{i\mathbf{k}\cdot\mathbf{r}}$；
    \item 能量色散关系为抛物线型 $E(\mathbf{k}) = \hbar^2k^2/2m$；
    \item 引入\textbf{费米能} $E_F$、\textbf{费米波矢} $k_F$、\textbf{费米温度} $T_F = E_F/k_{\text{B}}$ 和\textbf{费米面}等重要概念。
\end{enumerate}

\textbf{成功应用}：成功解释了电子的热容（$C_e = \gamma T$）、泡利顺磁性、热电子发射、接触电势差等经典理论无法说明的实验事实。

% --------------------------------------------------------
\item \textbf{请说明石墨的结构特点。}

石墨为\textbf{层状结构}：
\begin{enumerate}[nosep, label=(\arabic*)]
    \item \textbf{层内}：每个碳原子以 $\text{sp}^2$ 杂化与同一层内的三个碳原子形成强 $\sigma$ 共价键，构成正六角形蜂巢状平面网络。C--C键长约 $1.42$\AA，键能很高。
    \item \textbf{层间}：层与层之间靠范德瓦尔斯力结合，间距约 $3.35$\AA，结合力很弱，因此石墨易沿层间滑移，可用作润滑剂。
    \item \textbf{电子结构}：每个碳原子还有一个垂直于层面的 $\text{p}_z$ 轨道电子，形成离域的大 $\pi$ 键，使石墨在层面方向具有良好的导电性和金属光泽（半金属性）。
\end{enumerate}

石墨是\textbf{各向异性材料}的典型代表：层内硬度高、导电性好；层间软弱、易剥离。石墨烯（单层石墨）就是从石墨中剥离出的单原子层二维材料。

\end{enumerate}
